\documentclass[11pt]{article}

\usepackage[margin=1in]{geometry}
\usepackage{amsmath,amssymb,amsthm,mathtools}
\usepackage{hyperref}
\usepackage{enumitem}
\usepackage{booktabs}
\usepackage{microtype}

\raggedbottom
\allowdisplaybreaks[2]
\setlength{\parindent}{1.5em}
\setlength{\parskip}{0pt}
\setlist{topsep=3pt,itemsep=2pt,parsep=0pt,partopsep=0pt}

\hypersetup{colorlinks=true,linkcolor=blue,citecolor=blue,urlcolor=blue}

\newtheoremstyle{tightplain}{6pt}{6pt}{\itshape}{}{\bfseries}{.}{.5em}{}
\newtheoremstyle{tightdef}{6pt}{6pt}{}{}{\bfseries}{.}{.5em}{}
\theoremstyle{tightplain}
\newtheorem{theorem}{Theorem}[section]
\newtheorem{proposition}[theorem]{Proposition}
\newtheorem{lemma}[theorem]{Lemma}
\newtheorem{corollary}[theorem]{Corollary}
\newtheorem{openstep}[theorem]{Deferred step}
\theoremstyle{tightdef}
\newtheorem{assumption}[theorem]{Assumption}
\newtheorem{definition}[theorem]{Definition}
\newtheorem{remark}[theorem]{Remark}

\DeclareMathOperator{\tr}{tr}
\DeclareMathOperator{\diag}{diag}
\DeclareMathOperator{\rank}{rank}
\DeclareMathOperator{\op}{op}
\DeclareMathOperator{\KL}{KL}

\newcommand{\E}{\mathbb E}
\newcommand{\Prob}{\mathbb P}
\newcommand{\R}{\mathbb R}
\newcommand{\norm}[1]{\left\lVert #1 \right\rVert}
\newcommand{\abs}[1]{\left\lvert #1 \right\rvert}
\newcommand{\one}{\mathbf 1}
\newcommand{\F}{\mathrm F}
\newcommand{\Shat}{\widehat\Sigma}
\newcommand{\Phat}{\widehat P}
\newcommand{\Sbar}{\bar\Sigma}
\newcommand{\dsub}{\mathfrak d_{\mathrm{sub}}}
\newcommand{\Gcal}{\mathcal G}
\newcommand{\Ecal}{\mathcal E}
\newcommand{\Hcal}{\mathcal H}

\title{Adjacent-Window Fluctuations of\\
Rolling-Covariance Spectral Projectors\\
under a Static Gaussian Null}
\date{\today}

\begin{document}
\maketitle

\begin{abstract}
We study how the rank-$r$ spectral projector of a rolling-window sample covariance
matrix changes between two consecutive, overlapping windows when the data are
i.i.d.\ Gaussian. The two windows share all but one observation each, so the
covariance difference has rank two; we quantify the resulting squared projector
increment $\E\norm{\Phat_t-\Phat_{t-1}}_F^2$ and show it scales as $r/W^2$, with an
explicit constant. The self-contained core proved here consists of a conditional
exchangeability identity for the two projectors given the shared observations, an
exact conditional sign-randomization principle (which yields a finite-sample test for
oriented contrasts), exact Gaussian formulas for the linearized projector variance
both for a fixed population and conditionally around the random overlap covariance,
and the algebra giving the $r/W^2$ rate. The perturbation, concentration,
$\chi^2$-tail, and Gaussian-chaos facts used along the way are standard and are cited
at the point of use rather than restated; with those citations in place the $r/W^2$
rate (Theorem~\ref{thm:conditionalbenchmark}) holds outright. Only two items are
specific to this work and remain to be written: a Lipschitz-continuity step for the
conditional variance coefficient, needed for the sharp constant
(Theorem~\ref{thm:conditionalconstant}), and the local-geometry details of a matching
minimax lower bound.
\end{abstract}

\section{Conventions on proof status}

We use two labels. A \emph{Proposition}, \emph{Lemma}, \emph{Corollary}, or
\emph{Theorem} is proved here in full from elementary conditioning, Gaussian moment
identities, linear algebra, and standard results cited inline where they are used. A
\emph{Deferred step} is specific to this work and is presently only sketched; these
are the items a final version still has to write out.

The results proved here are: the rank-two structure and conditional exchangeability
(\S\ref{sec:exch}); exact conditional sign randomization and the resulting
finite-sample test (\S\ref{sec:sign}); the exact fixed-population linearized variance
(\S\ref{sec:fixedpop}); the exact conditional linearized variance around the random
overlap covariance, the overlap-separation lemma, and the order-$r$ control of the
variance coefficient (\S\ref{sec:condvar}); and the $r/W^2$ rate together with the
conditional constant-level expansion (\S\ref{sec:rate}). Standard perturbation,
covariance-concentration, $\chi^2$-tail, and Gaussian-chaos facts are invoked by
citation where needed and are not restated as separate claims. After this accounting,
the rate in Theorem~\ref{thm:conditionalbenchmark} is unconditional; the sharp
constant in Theorem~\ref{thm:conditionalconstant} rests on a single continuity step
(\S\ref{sec:deferred}, Deferred step~\ref{def:corederivative}), and the matching lower
bound on a second (Deferred step~\ref{def:rollinglower}).

\section{Setup and assumptions}

For a window width $W$ let $I_t=\{t-W+1,\ldots,t\}$ and define the rolling sample
covariance
\[
\widehat\Sigma_t=\frac1W\sum_{s\in I_t}r_s r_s^\top .
\]
For a symmetric matrix $A$ whose $r$-th and $(r+1)$-st eigenvalues differ, let $P(A)$
be its top-$r$ spectral projector; on the measure-zero set of eigenvalue ties we fix
an arbitrary measurable tie-breaking rule. Write $\Phat_t=P(\widehat\Sigma_t)$.

\begin{assumption}[Static Gaussian null]\label{ass:gaussian}
On $I_t\cup I_{t-1}$ the vectors $r_s$ are independent and identically distributed
$N(0,\Sigma)$.
\end{assumption}

\begin{assumption}[Bulk-plus-spikes spectrum]\label{ass:bulkspike}
The rank $r$ is fixed. With eigenvalues ordered as
$\lambda_1(\Sigma)\ge\cdots\ge\lambda_n(\Sigma)$, there are constants
$0<m<M<\infty$ and $0<c_\lambda\le C_\lambda<\infty$, independent of $n$ and $W$,
with
\[
m\le \lambda_n(\Sigma)\le\cdots\le\lambda_{r+1}(\Sigma)\le M,
\qquad
c_\lambda n\le \lambda_r(\Sigma)\le\cdots\le\lambda_1(\Sigma)\le C_\lambda n .
\]
\end{assumption}

\begin{assumption}[Eigengap]\label{ass:gap}
The gap $\delta=\lambda_r(\Sigma)-\lambda_{r+1}(\Sigma)$ satisfies $\delta\ge c_\delta n$
for a constant $c_\delta>0$.
\end{assumption}

\begin{assumption}[Asymptotic regime]\label{ass:asymptotic}
The rank $r$ is fixed, $n$ may grow, and $W=W_n\to\infty$. Wherever exponential
sample-covariance concentration is invoked we assume $W\ge C\log n$ and $W\ge W_0$,
with constants depending only on the spectral constants above and on $r$.
\end{assumption}

For a Hilbert-space-valued square-integrable random variable $X$ and a
sub-$\sigma$-algebra $\Gcal$, write the conditional variance
\[
V_F(X\mid\Gcal)=\E\bigl[\norm{X-\E[X\mid\Gcal]}_F^2\bigm|\Gcal\bigr].
\]

\begin{lemma}[Conditional i.i.d.\ variance identity]\label{lem:iidvar}
If $X_1,X_2$ are conditionally i.i.d.\ and square-integrable given $\Gcal$, then
\[
\E\bigl[\norm{X_1-X_2}_F^2\bigm|\Gcal\bigr]=2\,V_F(X_1\mid\Gcal).
\]
\end{lemma}

\begin{proof}
Let $\mu=\E[X_1\mid\Gcal]=\E[X_2\mid\Gcal]$. Conditional independence gives
$\E[\langle X_1-\mu,X_2-\mu\rangle\mid\Gcal]=0$, so
$\E[\norm{X_1-X_2}_F^2\mid\Gcal]=\E[\norm{X_1-\mu}_F^2\mid\Gcal]
+\E[\norm{X_2-\mu}_F^2\mid\Gcal]=2V_F(X_1\mid\Gcal)$.
\end{proof}

\section{Results proved here}

\subsection{Rank-two structure and conditional exchangeability}\label{sec:exch}

For adjacent windows,
\[
\widehat\Sigma_t-\widehat\Sigma_{t-1}
=\frac1W\bigl(r_t r_t^\top-r_{t-W}r_{t-W}^\top\bigr),
\]
so the covariance difference has rank at most two: the windows differ only through
the \emph{entering} observation $r_t$ and the \emph{leaving} observation $r_{t-W}$.
Their common part is the \emph{overlap covariance}
\[
C_t:=\frac1W\sum_{s=t-W+1}^{t-1}r_s r_s^\top,
\qquad
\alpha_W:=\frac{W-1}{W},
\]
which satisfies $\E C_t=\alpha_W\Sigma$ and
\[
\widehat\Sigma_t=C_t+\frac1W r_t r_t^\top,
\qquad
\widehat\Sigma_{t-1}=C_t+\frac1W r_{t-W}r_{t-W}^\top.
\]
Let $\phi_{C_t}(v)$ denote the top-$r$ projector of $C_t+W^{-1}vv^\top$, and set
$\Gcal_t=\sigma(r_{t-W+1},\ldots,r_{t-1})$.

\begin{proposition}[Conditional exchangeability]\label{prop:exchangeability}
Under Assumption~\ref{ass:gaussian}, conditional on $\Gcal_t$ the projectors
\[
\Phat_t=\phi_{C_t}(r_t),
\qquad
\Phat_{t-1}=\phi_{C_t}(r_{t-W})
\]
are i.i.d. Consequently
\[
\E[\Phat_t-\Phat_{t-1}\mid\Gcal_t]=0,
\qquad
\E\norm{\Phat_t-\Phat_{t-1}}_F^2=2\,\E\bigl[V_F(\Phat_t\mid\Gcal_t)\bigr].
\]
\end{proposition}

\begin{proof}
Given $\Gcal_t$ the overlap covariance $C_t$ is deterministic and the two boundary
observations $r_t,r_{t-W}$ are i.i.d.\ $N(0,\Sigma)$, independent of $C_t$. Hence
$\Phat_t$ and $\Phat_{t-1}$ are conditionally i.i.d. The mean-zero statement is then
immediate, and the second-moment statement is Lemma~\ref{lem:iidvar} followed by
taking expectations.
\end{proof}

\subsection{Sign randomization and a finite-sample test}\label{sec:sign}

\begin{theorem}[Exact conditional sign randomization]\label{thm:exactrandomization}
Under Assumption~\ref{ass:gaussian}, let $\mathcal A$ be any measurable
Hilbert-space-valued functional of two rank-$r$ projectors with
$\mathcal A(Q_1,Q_2)=-\mathcal A(Q_2,Q_1)$. Then, conditional on $\Gcal_t$,
\[
\mathcal A(\Phat_t,\Phat_{t-1})\stackrel{d}=-\mathcal A(\Phat_t,\Phat_{t-1}).
\]
This applies to oriented statistics such as $\Phat_t-\Phat_{t-1}$ and its linear
contrasts. It says nothing about the nonnegative magnitude
$\norm{\Phat_t-\Phat_{t-1}}_F$, which is an even functional.
\end{theorem}

\begin{proof}
Swapping the two conditionally i.i.d.\ boundary observations $r_t$ and $r_{t-W}$
preserves the conditional law and exchanges the two projectors; antisymmetry of
$\mathcal A$ then negates the statistic.
\end{proof}

\begin{corollary}[Finite-sample sign test for oriented contrasts]\label{cor:signtest}
Under Assumption~\ref{ass:gaussian}, fix a deterministic symmetric matrix $A$ or a
$\Gcal_t$-measurable symmetric matrix $A_t$, and set
$T_t=\langle A_t,\Phat_t-\Phat_{t-1}\rangle_F$. Then $T_t\stackrel d=-T_t$
conditional on $\Gcal_t$. Hence if $\Prob(T_t=0\mid\Gcal_t)=0$ the conditional signs
of $T_t$ are exactly balanced under the null; in the presence of an atom at zero,
treating zeros as non-rejections gives a conservative test.
\end{corollary}

\begin{proof}
Apply Theorem~\ref{thm:exactrandomization} to
$\mathcal A(Q_1,Q_2)=\langle A_t,Q_1-Q_2\rangle_F$, conditioning on $\Gcal_t$ when
$A_t$ is random.
\end{proof}

\subsection{Exact linearized variance: fixed population}\label{sec:fixedpop}

\begin{definition}[Linearized-projector variance coefficient]\label{def:dsub}
With $\Sigma=\sum_{k=1}^n\lambda_k u_k u_k^\top$, $\lambda_1\ge\cdots\ge\lambda_n$ and
$\lambda_r>\lambda_{r+1}$, set
\[
\dsub(\Sigma,r)=2\sum_{i\le r<j}\frac{\lambda_i\lambda_j}{(\lambda_i-\lambda_j)^2}.
\]
\end{definition}

\begin{proposition}[Exact fixed-population linearized variance]\label{prop:dsub}
Let $P=P(\Sigma)$ and let
\[
\mathcal L_\Sigma(E)=
\sum_{i\le r<j}
\frac{u_j u_j^\top E\, u_i u_i^\top+u_i u_i^\top E\, u_j u_j^\top}{\lambda_i-\lambda_j}
\]
be the first-order projector perturbation. Then
\[
\norm{\mathcal L_\Sigma(E)}_F^2
=2\sum_{i\le r<j}\frac{\abs{u_j^\top E u_i}^2}{(\lambda_i-\lambda_j)^2}.
\]
If $\widehat\Sigma=W^{-1}\sum_{s=1}^W r_s r_s^\top$ with i.i.d.\
$r_s\sim N(0,\Sigma)$, then
\[
\E\norm{\mathcal L_\Sigma(\widehat\Sigma-\Sigma)}_F^2=\frac{\dsub(\Sigma,r)}{W}.
\]
If $H=W^{-1}(r_+r_+^\top-r_-r_-^\top)$ with $r_+,r_-\stackrel{\mathrm{iid}}\sim N(0,\Sigma)$, then
\[
\E\norm{\mathcal L_\Sigma(H)}_F^2=\frac{2\,\dsub(\Sigma,r)}{W^2}.
\]
\end{proposition}

\begin{proof}
The Frobenius identity follows from the orthogonality of the matrix units
$u_j u_i^\top$ in the reduced-resolvent expression. For Gaussian data and $i\ne j$,
\[
\E\abs{u_j^\top(\widehat\Sigma-\Sigma)u_i}^2=\frac{\lambda_i\lambda_j}{W},
\qquad
\E\abs{u_j^\top H u_i}^2=\frac{2\lambda_i\lambda_j}{W^2},
\]
since $u_i^\top r_s$ and $u_j^\top r_s$ are independent mean-zero Gaussians with
variances $\lambda_i,\lambda_j$. Summing over $i\le r<j$ gives both claims.
\end{proof}

\begin{corollary}[Bulk-plus-spikes scale]\label{cor:firstorderscale}
Under Assumptions~\ref{ass:bulkspike}--\ref{ass:gap}, $\dsub(\Sigma,r)=\Theta(r)$.
Hence the single-window linearized variance scales as $r/W$ and the adjacent-window
linearized variance as $r/W^2$.
\end{corollary}

\begin{proof}
For $i\le r<j$ one has $\lambda_i\asymp n$, $\lambda_j\asymp 1$, and
$\lambda_i-\lambda_j\asymp n$, so each summand is $\asymp 1/n$. There are
$r(n-r)$ summands, giving $\Theta(r)$ for fixed $r$.
\end{proof}

\subsection{Exact conditional linearized variance: random overlap}\label{sec:condvar}

Write $C_t=\sum_{k=1}^n\widetilde\lambda_k\widetilde u_k\widetilde u_k^\top$ with
$\widetilde\lambda_1\ge\cdots\ge\widetilde\lambda_n$. On the event
$\widetilde\lambda_r>\widetilde\lambda_{r+1}$ define $\mathcal L_{C_t}$ as in
Proposition~\ref{prop:dsub} but with the eigendata of $C_t$, and set
\[
a_i=\widetilde u_i^\top\Sigma\widetilde u_i,
\qquad
b_j=\widetilde u_j^\top\Sigma\widetilde u_j,
\qquad
c_{ij}=\widetilde u_i^\top\Sigma\widetilde u_j .
\]

\begin{proposition}[Exact conditional linearized variance]\label{prop:conditionallinearized}
Under Assumption~\ref{ass:gaussian}, on $\widetilde\lambda_r>\widetilde\lambda_{r+1}$,
with $H_t=W^{-1}(r_t r_t^\top-r_{t-W}r_{t-W}^\top)$,
\[
\E\bigl[\norm{\mathcal L_{C_t}(H_t)}_F^2\bigm|\Gcal_t\bigr]
=\frac4{W^2}\sum_{i\le r<j}\frac{a_i b_j+c_{ij}^2}{(\widetilde\lambda_i-\widetilde\lambda_j)^2}.
\]
\end{proposition}

\begin{proof}
For one boundary term $\Delta=xx^\top/W$ with $x\sim N(0,\Sigma)$,
\[
\norm{\mathcal L_{C_t}(\Delta)}_F^2
=\frac2{W^2}\sum_{i\le r<j}
\frac{\abs{\widetilde u_j^\top x}^2\abs{\widetilde u_i^\top x}^2}
{(\widetilde\lambda_i-\widetilde\lambda_j)^2}.
\]
With $g_i=\widetilde u_i^\top x$, Isserlis' formula gives
$\E[g_i^2 g_j^2\mid\Gcal_t]=a_i b_j+2c_{ij}^2$, while
$\norm{\E[\mathcal L_{C_t}(\Delta)\mid\Gcal_t]}_F^2
=\tfrac2{W^2}\sum_{i\le r<j}c_{ij}^2/(\widetilde\lambda_i-\widetilde\lambda_j)^2$.
Subtracting gives the conditional variance of one term; the two terms are
conditionally i.i.d., so the second moment of their difference is twice that, which
is the displayed expression.
\end{proof}

Define the conditional variance coefficient
\[
Q(C_t,\Sigma)
=4\sum_{i\le r<j}\frac{a_i b_j+c_{ij}^2}{(\widetilde\lambda_i-\widetilde\lambda_j)^2},
\qquad\text{so}\qquad
\E\bigl[\norm{\mathcal L_{C_t}(H_t)}_F^2\bigm|\Gcal_t\bigr]=\frac{Q(C_t,\Sigma)}{W^2},
\]
and the concentration event
\[
\Ecal_t=\bigl\{\norm{C_t-\alpha_W\Sigma}_{\op}\le \eta\,\alpha_W n\bigr\},
\]
where $\eta>0$ is small relative to $c_\lambda$ and $c_\delta$.

\begin{lemma}[Overlap-covariance separation]\label{inp:coresep}
Under Assumptions~\ref{ass:gaussian}--\ref{ass:gap}, for $\eta>0$ small there are
$c,C>0$ with $\Prob(\Ecal_t^c)\le Ce^{-cW}$ whenever $W\ge C\log n$ and $W\ge W_0$.
\end{lemma}

\begin{proof}
Write $C_t=\alpha_W S_t$ with $S_t=(W-1)^{-1}\sum_{s=t-W+1}^{t-1}r_s r_s^\top$, the
sample covariance of $W-1$ i.i.d.\ $N(0,\Sigma)$ vectors. Then
$\Ecal_t^c=\{\norm{S_t-\Sigma}_{\op}>\eta n\}$, and the bound is the exponential
operator-norm tail of the Gaussian sample-covariance concentration inequality
(Koltchinskii and Lounici~\cite{kl}; Vershynin~\cite{vershynin}, Ch.~4--6) applied
with $N=W-1$ and effective rank $\mathbf r(\Sigma)=O(1)$ under
Assumption~\ref{ass:bulkspike}.
\end{proof}

\begin{lemma}[Order-$r$ control of the variance coefficient]\label{lem:sharedcorelinearized}
Under Assumptions~\ref{ass:gaussian}--\ref{ass:gap}, on $\Ecal_t$, and with $W\ge W_0$
so that $\alpha_W$ is bounded away from zero,
\[
c r\le Q(C_t,\Sigma)\le C r .
\]
Consequently, by Lemma~\ref{inp:coresep},
$\E\bigl[\norm{\mathcal L_{C_t}(H_t)}_F^2\,\one_{\Ecal_t}\bigr]\asymp r/W^2$.
\end{lemma}

\begin{proof}
On $\Ecal_t$, Weyl's inequality gives
$\widetilde\lambda_r-\widetilde\lambda_{r+1}\ge\alpha_W\delta-2\eta\alpha_W n
\ge\tfrac12\alpha_W\delta\asymp n$, so $\abs{\widetilde\lambda_i-\widetilde\lambda_j}\gtrsim n$
for $i\le r<j$; together with $0\le\widetilde\lambda_j\le\widetilde\lambda_i\lesssim n$
the denominators are of order $n^2$.

By Cauchy--Schwarz for the form $\langle\cdot,\cdot\rangle_\Sigma$, $c_{ij}^2\le a_i b_j$,
so $\sum_{i\le r<j}(a_i b_j+c_{ij}^2)\le 2(\sum_{i\le r}a_i)(\sum_{j>r}b_j)$. Since
$\Sigma$ has $r$ spikes of size $O(n)$ and bulk trace $O(n)$,
$\sum_{i\le r}a_i\lesssim rn$ and $\sum_{j>r}b_j\lesssim n$, giving the upper bound.

For the lower bound, on $\Ecal_t$ one has $\widetilde\lambda_i\gtrsim n$ for $i\le r$, and
since $\widetilde u_i$ is an eigenvector of $C_t$,
$\alpha_W a_i=\widetilde u_i^\top(\alpha_W\Sigma)\widetilde u_i
\ge\widetilde\lambda_i-\norm{C_t-\alpha_W\Sigma}_{\op}$, so $a_i\gtrsim n$ for $i\le r$
after fixing $\eta$ small. By Ky Fan,
$\sum_{j>r}b_j=\tr((I-P(C_t))\Sigma)\ge\sum_{k=r+1}^n\lambda_k(\Sigma)\gtrsim n$.
Dropping $c_{ij}^2\ge0$ and bounding denominators by $Cn^2$ yields the lower bound.
The expectation statement combines this with Lemma~\ref{inp:coresep}.
\end{proof}

\subsection{Rate and constant}\label{sec:rate}

\begin{theorem}[Adjacent-window rate]\label{thm:conditionalbenchmark}
Assume Assumptions~\ref{ass:gaussian}--\ref{ass:gap} with $r$ fixed. Then, for
$W\ge C\log n$ and $W\ge W_0$,
\[
c\,\frac r{W^2}\le \E\norm{\Phat_t-\Phat_{t-1}}_F^2\le C\,\frac r{W^2}+Cr e^{-cW}.
\]
\end{theorem}

\begin{proof}
Let $\gamma_t=\widetilde\lambda_r-\widetilde\lambda_{r+1}$ and
$\Delta_t^+=W^{-1}r_t r_t^\top$, $\Delta_t^-=W^{-1}r_{t-W}r_{t-W}^\top$,
$H_t=\Delta_t^+-\Delta_t^-$. Throughout we use the second-order perturbation
expansion of a top-$r$ spectral projector with a gap-controlled quadratic remainder:
for symmetric $A$ with gap $\gamma=\mu_r-\mu_{r+1}>0$ and
$\norm E_{\op}\le c_0\gamma$,
\[
\norm{P(A+E)-P(A)-\mathcal L_A(E)}_F\le C_r\norm E_{\op}^2/\gamma^2,
\]
with $c_0\in(0,1/4]$ and $C_r<\infty$ (Kato~\cite{kato}, Ch.~II; Stewart and
Sun~\cite{stewartsun}, Ch.~V).

Fix $\rho\in(0,c_0]$ and set
$\Hcal_t=\Ecal_t\cap\{\norm{\Delta_t^+}_{\op}\vee\norm{\Delta_t^-}_{\op}\le\rho\gamma_t\}$.
Lemma~\ref{inp:coresep} controls $\Ecal_t^c$, and the boundary norms are controlled by
a weighted-$\chi^2$ tail: since $\norm{\Delta_t^\pm}_{\op}\le\norm{r_\pm}_2^2/W$ and
$\E\norm{r_\pm}_2^2=\tr\Sigma\asymp n$, one has
$\Prob\{\norm{r_\pm}_2^2>\kappa nW\}\le C_\kappa e^{-c_\kappa W}$ for each fixed
$\kappa>0$ (Laurent and Massart~\cite{laurentmassart}, Lemma~1; equivalently the
Hanson--Wright inequality, Rudelson and Vershynin~\cite{rudelsonvershynin}). Hence
$\Prob(\Hcal_t^c)\le Ce^{-cW}$.

On $\Hcal_t$, the projector expansion applied at $C_t$ to each boundary term gives
$\Phat_t-\Phat_{t-1}=L_t+R_t^\Delta$ with $L_t=\mathcal L_{C_t}(H_t)$,
$R_t^\Delta=R_t^+-R_t^-$, and $\norm{R_t^\pm}_F\le C_r\norm{\Delta_t^\pm}_{\op}^2/\gamma_t^2$.
Since $\gamma_t\asymp n$ on $\Hcal_t$, the Gaussian eighth-moment bound yields
$\E[\norm{R_t^\Delta}_F^2\one_{\Hcal_t}]\le C_r W^{-4}$, while
Lemma~\ref{lem:sharedcorelinearized} gives $\E[\norm{L_t}_F^2\one_{\Ecal_t}]\asymp r/W^2$.
Conditionally on $\Gcal_t$, $L_t$ is a degree-two Gaussian chaos, whose $L^4$ and
$L^2$ norms are equivalent by hypercontractivity (Janson~\cite{janson}, Ch.~5--6);
with the second moment $\asymp r/W^2$ this gives
$\E[\norm{L_t}_F^4\mid\Gcal_t]\le Cr^2/W^4$ on $\Ecal_t$, so the loss on
$\Ecal_t\cap\Hcal_t^c$ is exponentially small and $\E[\norm{L_t}_F^2\one_{\Hcal_t}]\asymp r/W^2$.

For the upper bound use
$\norm{L_t+R_t^\Delta}_F^2\le 2\norm{L_t}_F^2+2\norm{R_t^\Delta}_F^2$ on $\Hcal_t$ and
$\norm{\Phat_t-\Phat_{t-1}}_F^2\le 2r$ off $\Hcal_t$. For the lower bound,
$\norm{L_t+R_t^\Delta}_F^2\ge\norm{L_t}_F^2-2\abs{\langle L_t,R_t^\Delta\rangle}$ and
Cauchy--Schwarz give a cross term of order $\sqrt r\,W^{-3}=o(r/W^2)$.
\end{proof}

\begin{theorem}[Constant-level expansion; conditional on Deferred step~\ref{def:corederivative}]\label{thm:conditionalconstant}
Assume Assumptions~\ref{ass:gaussian}--\ref{ass:gap} with $r$ fixed, $W\to\infty$,
$W\ge C\log n$. Granting Deferred step~\ref{def:corederivative},
\[
\E\norm{\Phat_t-\Phat_{t-1}}_F^2
=\frac{2}{\alpha_W^2}\frac{\dsub(\Sigma,r)}{W^2}+o(W^{-2})
=\frac{2\dsub(\Sigma,r)}{W^2}+o(W^{-2}),
\]
the last equality because $\alpha_W\to1$ and $\dsub(\Sigma,r)=O(1)$ for fixed $r$.
\end{theorem}

\begin{proof}
Use the event $\Hcal_t$ and decomposition $\Phat_t-\Phat_{t-1}=L_t+R_t^\Delta$ of
Theorem~\ref{thm:conditionalbenchmark}, with
$\E[\norm{R_t^\Delta}_F^2\one_{\Hcal_t}]\le CW^{-4}$ and $\Prob(\Hcal_t^c)\le Ce^{-cW}$.
By Proposition~\ref{prop:conditionallinearized},
$\E[\norm{L_t}_F^2\mid\Gcal_t]=Q(C_t,\Sigma)/W^2$ on the gap event. Deferred
step~\ref{def:corederivative} with Lemma~\ref{inp:coresep} gives
$\E\norm{L_t}_F^2=W^{-2}\bigl(2\alpha_W^{-2}\dsub(\Sigma,r)+o(1)\bigr)+O(e^{-cW})$; the
restriction loss to $\Hcal_t$ is exponentially small by the conditional
fourth-moment bound (hypercontractivity; Janson~\cite{janson}, Ch.~5--6). Expanding
$\norm{L_t+R_t^\Delta}_F^2$, the remainder square is $O(W^{-4})$ and the cross term
$O(W^{-3})$, both $o(W^{-2})$.
\end{proof}

\begin{remark}[Consistency check on the constant]
Replacing the random overlap covariance by its mean $C_t=\alpha_W\Sigma$ gives
$\widetilde u_i=u_i$, $a_i=\lambda_i$, $b_j=\lambda_j$, $c_{ij}=0$, hence
$Q(\alpha_W\Sigma,\Sigma)=4\sum_{i\le r<j}\lambda_i\lambda_j/(\alpha_W^2(\lambda_i-\lambda_j)^2)
=2\alpha_W^{-2}\dsub(\Sigma,r)$, matching Theorem~\ref{thm:conditionalconstant}.
Deferred step~\ref{def:corederivative} is precisely the statement that this
replacement is asymptotically valid.
\end{remark}

\section{Deferred steps (specific to this work)}\label{sec:deferred}

These two statements are not off-the-shelf theorems: their ingredients are standard,
but the assembled quantitative claims are particular to this problem and still have to
be written out.

\begin{openstep}[Continuity of the variance coefficient]\label{def:corederivative}
Under Assumptions~\ref{ass:gaussian}--\ref{ass:gap} with $r$ fixed, $W\to\infty$,
$W\ge C\log n$,
\[
\E\left[\abs[\Big]{Q(C_t,\Sigma)-\frac{2}{\alpha_W^2}\dsub(\Sigma,r)}\,\one_{\Ecal_t}\right]
\le C\,\frac{\E\norm{C_t-\alpha_W\Sigma}_{\op}}{n}=O(W^{-1/2})=o(1).
\]
\emph{The load-bearing step.} The claim is that the basis-free map
$A\mapsto Q(A,\Sigma)=\E\norm{\mathcal L_A(xx^\top-yy^\top)}_F^2$, with
$x,y\stackrel{\mathrm{iid}}\sim N(0,\Sigma)$, is Lipschitz at $A_0=\alpha_W\Sigma$ with
modulus $C/n$ on an operator-norm ball of radius $O(n)$. The intended proof combines
the Riesz-projector representation of $\mathcal L_A$ and the resolvent identity (the
perturbation machinery of Kato~\cite{kato}, Ch.~II) with fixed-rank bulk-plus-spikes
moment bounds (Koltchinskii and Lounici~\cite{kl}). This is the one step that should
be written in full before submission.
\end{openstep}

\begin{openstep}[Minimax lower bound for the projector increment]\label{def:rollinglower}
Fix $r$ with $n\ge 2r+1$. There is a Gaussian adjacent-window subclass meeting the
bulk-plus-spikes and eigengap conditions uniformly such that, with
$D_t=P(\bar\Sigma_t)-P(\bar\Sigma_{t-1})$,
\[
\inf_{\widehat D}\sup\,\E\norm{\widehat D-D_t}_F^2\ge c\,\frac r{W^2}.
\]
\emph{Status.} The intended proof uses the local chart
$U_\Theta=\binom{I_r}{\Theta}(I_r+\Theta^\top\Theta)^{-1/2}$,
$P_\Theta=U_\Theta U_\Theta^\top$, $\Sigma_\Theta=I_n+nP_\Theta$, with the overlap
covariance set to $\Sigma_0$ and the two boundary covariances to
$\Sigma_{\pm\Theta}$. The standard machinery is Assouad's lemma together with the
Gaussian KL formula (Tsybakov~\cite{tsybakov}, \S2.7); the work specific to this paper
is the local projector-geometry equivalence, the first-order expansion of
$D_t(\Theta)$, the exact KL between adjacent hypercube vertices, and the Assouad
reduction.
\end{openstep}

\section{Combined statement}

\begin{theorem}[Adjacent-window projector increment]\label{thm:final}
Assume Assumptions~\ref{ass:gaussian}--\ref{ass:asymptotic}. The rate
\[
\E\norm{\Phat_t-\Phat_{t-1}}_F^2\asymp \frac r{W^2}
\]
holds unconditionally by Theorem~\ref{thm:conditionalbenchmark}. Granting the single
continuity step (Deferred step~\ref{def:corederivative}), the sharp constant is
\[
\E\norm{\Phat_t-\Phat_{t-1}}_F^2=\frac{2\dsub(\Sigma,r)}{W^2}+o(W^{-2}),
\qquad \dsub(\Sigma,r)=\Theta(r).
\]
\end{theorem}

\begin{remark}
Under the static null the conditional mean of the increment is zero
(Proposition~\ref{prop:exchangeability}), so the $r/W^2$ scale is the baseline
sampling-fluctuation level of consecutive projectors, against which a genuine change
in the leading subspace must be detected. With the standard results cited inline, the
rate is settled; the remaining gaps are the constant
(Deferred step~\ref{def:corederivative}) and the matching lower bound
(Deferred step~\ref{def:rollinglower}). Distributional and concentration properties of
the spectral projectors of a single sample covariance, the closest prior object, are
studied by Koltchinskii and Lounici~\cite{klproj}.
\end{remark}

\begin{thebibliography}{9}

\bibitem{kato} T.~Kato. \emph{Perturbation Theory for Linear Operators}. Classics in
Mathematics. Springer, Berlin, 1995. Reprint of the 1980 second edition.

\bibitem{stewartsun} G.~W.~Stewart and J.-G.~Sun. \emph{Matrix Perturbation Theory}.
Academic Press, Boston, 1990.

\bibitem{kl} V.~Koltchinskii and K.~Lounici. Concentration inequalities and moment
bounds for sample covariance operators. \emph{Bernoulli}, 23(1):110--133, 2017.

\bibitem{klproj} V.~Koltchinskii and K.~Lounici. Normal approximation and
concentration of spectral projectors of sample covariance. \emph{Ann.\ Statist.},
45(1):121--157, 2017.

\bibitem{vershynin} R.~Vershynin. \emph{High-Dimensional Probability: An Introduction
with Applications in Data Science}. Cambridge Series in Statistical and Probabilistic
Mathematics. Cambridge University Press, Cambridge, 2018.

\bibitem{laurentmassart} B.~Laurent and P.~Massart. Adaptive estimation of a quadratic
functional by model selection. \emph{Ann.\ Statist.}, 28(5):1302--1338, 2000.

\bibitem{rudelsonvershynin} M.~Rudelson and R.~Vershynin. Hanson--Wright inequality and
sub-gaussian concentration. \emph{Electron.\ Commun.\ Probab.}, 18(82):1--9, 2013.

\bibitem{janson} S.~Janson. \emph{Gaussian Hilbert Spaces}. Cambridge Tracts in
Mathematics 129. Cambridge University Press, Cambridge, 1997.

\bibitem{tsybakov} A.~B.~Tsybakov. \emph{Introduction to Nonparametric Estimation}.
Springer Series in Statistics. Springer, New York, 2009.

\end{thebibliography}

\end{document}